\section{Open Questions}

This version leaves several questions intentionally open.

\begin{enumerate}[label=\textbf{Q\arabic*.}]
    \item \textbf{How does the third-order DGP map into the first-order CKM accounting VAR?}
    The simulation DGP is solved by third-order perturbation with pruning,
    because volatility shocks affect allocations only beyond certainty
    equivalence. The remaining methodological question is how the higher-order
    policy-function terms generated by that DGP map into the first-order VAR
    expectations used when recovering empirical CKM wedges. The local-projection
    and GIRF exercises should quantify this approximation gap rather than
    revisit whether the DGP itself must be third-order.

    \item \textbf{How should the investment wedge be timed?}
    This note places \(\tau_{x,t}\) at date \(t\) as a wedge on the expected
    return to capital. Some BCA normalizations place the investment wedge on the
    price of investment goods or index it at \(t+1\). The convention should be
    aligned with the CKM implementation before estimation.

    \item \textbf{Which wedges represent TFP uncertainty in practice?}
    The simulation exercise will show whether a TFP uncertainty shock is
    represented mainly as an investment wedge, a labor wedge, or a combination
    of canonical wedges in BCA space.

    \item \textbf{How will this connect to empirical uncertainty factors?}
    This note has one volatility state \(q_t\). The larger research project has
    empirical macro and financial uncertainty factors. The future SOE version
    will need a factor-loading structure that maps those empirical factors into
    volatility processes for canonical wedges.

    \item \textbf{How should SOE wedges enter?}
    In the small-open-economy version, the resource wedge should be decomposed
    into trade/resource absorption and possibly a UIP or country-premium wedge.
    This requires extending the closed-economy prototype without returning to
    NK price rigidity.
\end{enumerate}
